% PNP RESEARCH REPORT::FINAL-LINK2-REDTEAM::2026-10-08T01:02Z
\documentclass[11pt]{article}
\usepackage[T1]{fontenc}
\usepackage[margin=0.82in]{geometry}
\usepackage{amsmath,amssymb,amsthm,booktabs,longtable,array,xcolor,xurl,hyperref}
\hypersetup{colorlinks=true,linkcolor=blue!55!black,urlcolor=blue!55!black,citecolor=blue!55!black}
\setlength{\parskip}{0.28em}
\setlength{\parindent}{0pt}
\setlength{\emergencystretch}{3em}
\newtheorem{theorem}{Theorem}[section]
\newtheorem{lemma}[theorem]{Lemma}
\newtheorem{proposition}[theorem]{Proposition}
\newtheorem{corollary}[theorem]{Corollary}
\newcommand{\SAT}{\operatorname{SAT}}
\newcommand{\poly}{\operatorname{poly}}
\newcommand{\Exact}{\textsc{Proved/Exact}}
\newcommand{\Computational}{\textsc{Computational}}
\newcommand{\Open}{\textsc{Open}}
\newcommand{\Scoped}{\textsc{Scoped No-Go}}
\newcolumntype{L}[1]{>{\raggedright\arraybackslash}p{#1}}

\title{FINAL-LINK2 Red-Team Audit\\
Shared Acquisition, Delta Children, Clause-First DAGs,\\
Amortization, Typed Cycles, and Code-Dependent Portfolios}
\author{RED final-link run}
\date{8 October 2026, 01:02 UTC}

\begin{document}
\maketitle

\begin{abstract}
This report attacks the principal relaxations left by FINAL-LINK1.  The main
correction is positive but narrowing: a linear source-read floor does not
block the Williams region.  At the clean target, $M^{1/3}$ children of
$M^{2/3}$ bits have only $M$ total bits, so one shared scan and streamed
materialization are compatible with near-linear work and $M^{2/3}$ live
space.  Repeated full-parent scans cost $M^{4/3}$, but their necessity is not
an information-theoretic consequence of source size.

The audit proves restricted counterexamples to stronger shortcuts.  A small
delta can name a parent-scale residual; a clause-first cover can preserve an
arbitrary parent-scale payload; exact-residual DAG merging has an explicit
$2^d$-state family even after unit propagation; QLEN amortization is tight on
equal maximum-size children; a contracting typed cycle can still miss the
endpoint; and a parent-trace acquisition rule is destroyed by extensionally
irrelevant polynomial delays in the supplied SAT code.  A disjoint
pointer-chain oracle supplies a precise black-box direct-sum control.  None of
these is a general SAT lower bound.  A genuinely semantic, streamed,
acquisition-covariant BatchAcquire construction remains logically viable and
\Open.  Thus FINAL-LINK1 is not closed and P versus NP remains open.
\end{abstract}

\section{Claim card}

\begin{center}
\fbox{\begin{minipage}{0.92\textwidth}
\textbf{Result.}  FINAL-LINK2 does not prove or disprove the missing
successor theorem.  It rules out six tempting inferences:
\begin{enumerate}
\item a one-read lower bound rules out the $1/3,2/3$ cell;
\item a short delta handle is a small ordinary SAT instance;
\item clause-first disjunction plus exact residual equality guarantees a
small DAG;
\item aggregate query-length accounting automatically shares semantic SAT
work;
\item cycle product below one alone reaches the Williams endpoint; and
\item arbitrary code dependence permits cheap acquisition of a parent trace.
\end{enumerate}
The weakest survivor is an exact streamed typed successor whose common
physical-size Bellman ledger, aggregate decoder traffic, failed members, and
workspace all fit below the endpoint.  No current artifact constructs it on
all CNFs.
\end{minipage}}
\end{center}

All ``no-go'' statements below are restricted to the explicitly declared
representation or oracle model.  They must not be promoted to unrestricted
SAT lower bounds.  Conversely, the positive accounting observations are
possibility results, not constructions.

\section{Source reads do not by themselves block the target}

Let the parent have $M$ bits, let $q=M^{a+o(1)}$, and let every explicit child
have at most $L=M^{b+o(1)}$ bits.

\begin{lemma}[Equality-frontier traffic]
If $a+b=1$, then writing all children once costs at most $M^{1+o(1)}$ bits of
traffic.  A depth-first implementation can use one child buffer of
$M^{b+o(1)}$ bits.  Reading the parent once is also $M^{1+o(1)}$.  In
contrast, rescanning the full parent separately for every child costs
$M^{1+a+o(1)}$.
\end{lemma}

\begin{proof}
The aggregate child volume is $qL=M^{a+b+o(1)}=M^{1+o(1)}$.  Sequential
generation and consumption retain one child at a time.  One source scan adds
$M$ work; $q$ source scans add $qM=M^{1+a+o(1)}$.
\end{proof}

At $(a,b)=(1/3,2/3)$, the three relevant quantities are therefore
\[
 \text{one source scan}=M,qquad
 \text{all explicit child bits}=M,qquad
 \text{one root scan per child}=M^{4/3}.
\]
The first two fit the desired $M^{1+o(1)}$ acquisition budget.  Any argument
based only on ``the algorithm must inspect the input'' is consequently too
weak.  A useful negative theorem would have to force repeated access,
superlinear semantic computation, or an incompatible time-space tradeoff.

\subsection{A pivotal family calibrates the floor}

Define
\[
 U_n=(x_1)\wedge\cdots\wedge(x_n)\wedge
     (\neg x_1\vee\cdots\vee\neg x_n).
\]
$U_n$ is minimally unsatisfiable.  Replacing unit clause $(x_i)$ by
$(\neg x_i)$ makes the formula satisfiable without changing its clause width;
replacing the final all-negative clause by the all-positive clause also makes
it satisfiable without changing its width.  Hence every clause cell is
pivotal in a clause-oracle encoding, and an exact deterministic algorithm on
this input must inspect all $n+1$ clause cells.  The executable checked this
and deletion pivotality for $2\le n\le12$.

This is a real read lower bound and it still does not hurt the target.  In the
declared indexed bit encoding, $|U_n|=\Theta(n\log n)$, so even an
$\Omega(M)$ bit-level strengthening would remain admissible.  The right
conclusion is not ``reads are free,'' but ``one linear pass is affordable.''

\section{Shared BatchAcquire: what is and is not refuted}

\begin{proposition}[Batch composition is already the algorithm]
Suppose a uniform procedure constructs, on every $M$-bit CNF, an exact
counterfactually sufficient object $H$ in $M^{h+o(1)}$ time and the declared
space.  Suppose $H$ supports every recursive child access, combination, and
replay operation in a Bellman ledger with root budgets $(C,S)$.  Then the
procedure together with those consumers is a SAT algorithm with budgets
$(C,S)$; the name ``preprocessing'' gives no additional resource exemption.
\end{proposition}

The proposition is elementary but important.  BatchAcquire is not circular
merely because it is powerful; it is the exact missing constructive theorem.
It is circular only when one of its promised labels was itself obtained by a
parent-scale SAT call or an uncharged replay.

\subsection{A black-box direct-sum control}

Consider $q$ disjoint oracle paths.  Path $i$ begins at a known label and has
$R$ unknown next-pointer cells followed by one terminal bit.  Labels from
different paths are disjoint, and the output must contain all $q$ terminal
bits.

\begin{lemma}[Disjoint pointer-chain direct sum]
Every exact deterministic algorithm in this oracle model makes at least
$q(R+1)$ cell queries in the worst case.  Interleaving paths does not improve
the count.
\end{lemma}

\begin{proof}
The next address on a path is unavailable before its predecessor is queried,
so reaching one terminal needs all $R$ pointer queries.  Its terminal cell is
pivotal: if it is unqueried, flipping only that bit preserves the complete
observed transcript and changes a required output.  Disjoint paths share no
queried cell, so the lower bounds add.
\end{proof}

The experiment instantiated all $(q,R)$ with $q\in\{1,2,4,8\}$ and
$R\in\{2,4,8,16\}$ and obtained the exact count in every case.  This refutes
a representation-free claim that multiple calls to an arbitrary black-box
$A$ can always be batched.  It does \emph{not} refute a non-black-box sharing
theorem exploiting CNF semantics or the finite code of $A$.

\subsection{Why a one-run trace is not the required batch}

Let $B$ be any exact SAT decider.  For a constant $d$, let $A_d$ run $B$,
perform $n^d$ irrelevant steps, and then return $B$'s answer.  All $A_d$ are
extensionally identical SAT deciders.  Any acquisition contract requiring
reproduction of the full parent trace of the supplied code has exponent at
least $d$.  Since $d$ is arbitrary, that contract cannot have a common
endpoint bound.  Static code analysis may bypass a visible delay, so this is
a no-go for \emph{trace-defined} acquisition, not for every non-black-box
compiler.

The earlier project guard ``a post-pass cannot reduce already paid work'' is
therefore still active.  FINAL-LINK2 adds that trace cost is also non-covariant
under harmless implementation padding.

\section{Compressed delta children can hide a full child}

Let $P$ be any CNF on variables disjoint from
$C_w=(x_1\vee\cdots\vee x_w)$, and set $F=C_w\wedge P$.

\begin{lemma}[Payload inheritance]
For every clause-first term
\[
 \tau_i=\{x_1=0,\ldots,x_{i-1}=0,x_i=1\},
\]
the normalized residual $F\upharpoonright\tau_i$ is exactly $P$.  Thus all
children can merge, but their one remaining child has expanded size $|P|$.
If $|P|=\Theta(|F|)$, there is no power-size contraction.
\end{lemma}

\begin{proof}
$\tau_i$ satisfies $C_w$.  No assigned variable occurs in $P$, so restriction
does not alter $P$.
\end{proof}

The finite controls used equal-width, duplicate-free positive 3-CNF payloads.
At the largest tested point, the parent had 2,329 indexed bits, every one of
eight children normalized to the same 2,280-bit payload, while a crude delta
handle needed only 48 bits.  The handle exponent was about $0.499$ but the
expanded-child exponent was $0.997$.  Counting the former as the latter is a
type error.

\begin{lemma}[Honest delta-cycle scale]
Suppose an explicit $M$-bit source is replaced by a handle of
$M^{b+o(1)}$ bits, $0<b<1$, but consuming that handle as an ordinary SAT input
requires expanding a logical child of $\Theta(M)$ bits.  Relative to handle
size, the expansion edge has exponent at least $1/b$.  The round-trip scale
product is at least $b(1/b)=1$.
\end{lemma}

This does not forbid virtual children.  It forces the honest alternative:
keep the handle as a distinct typed problem, give a bit-level decoder, charge
every decoder call and retained source dependency, and prove a model-matched
recursive evaluator.  One may not inherit the ordinary-SAT bound at the
handle length while silently expanding a parent-scale CNF.

\section{Clause-first DAGs: an exact equality-merging control}

The payload lemma shows failure of contraction even when all routes merge.
The next family shows that exact residual equality need not merge the routes.

For $d\ge2$, introduce exact-one address gadgets
\[
 E_i=(a_i\vee b_i)\wedge(\neg a_i\vee\neg b_i),
 \qquad 1\le i\le d.
\]
For every address $\alpha\in\{0,1\}^d$, introduce $y_\alpha,u_\alpha,v_\alpha$
and clauses
\begin{align*}
 S_\alpha&=\left(\bigvee_{i:\alpha_i=0}a_i\right)
 \vee\left(\bigvee_{i:\alpha_i=1}\neg a_i\right)\vee y_\alpha,\\
 T_\alpha&=(\neg y_\alpha\vee u_\alpha\vee v_\alpha).
\end{align*}
Let $D_d$ be the conjunction of all $E_i,S_\alpha,T_\alpha$.  Apply the
two-term clause-first cover to $(a_i\vee b_i)$ in the order $i=1,\ldots,d$,
with unit propagation and syntactic normalization after each step.

\begin{theorem}[Clause-first equality-DAG gate]
After depth $h$, $D_d$ has exactly $2^h$ distinct normalized residuals, and
every branch extends to a model.  At depth $d$, every residual contains all
$2^d$ tag clauses, with exactly one $T_\alpha$ shortened to
$(u_\alpha\vee v_\alpha)$.  Therefore no two full-address residuals merge by
syntactic equality.
\end{theorem}

\begin{proof}
The two clause-first terms set $a_i=1,b_i=0$ or $a_i=0,b_i=1$ after the
exclusion clause propagates.  A prefix $p$ satisfies every selector whose
address disagrees with $p$ and retains precisely the selectors extending
$p$.  Different prefixes retain different indexed selector sets, proving
distinctness at depth $h$.  At a complete address $\alpha$, $S_\alpha$
becomes the unit $y_\alpha$ while every other selector is satisfied.  Unit
propagation shortens only $T_\alpha$.  Finally, set all $y_\alpha$ and all
$u_\alpha$ true to extend every address branch to a model.
\end{proof}

$D_d$ has $2d+2^{d+1}$ clauses and $2d+3\cdot2^d$ variables, so its indexed
length is $M=\Theta(2^d d)$.  The terminal DAG has $2^d=\Theta(M/\log M)$
distinct residuals and each retains $\Theta(M)$ indexed bits.  This rules out
the universal inference
\[
 \text{clause-first cover}+\text{UP}+\text{exact equality merging}
 \Longrightarrow (q,L)=(M^{1/3+o(1)},M^{2/3+o(1)}).
\]
It is deliberately not a SAT-hard family: all branches are satisfiable and a
semantic VALUE routine can bypass the DAG.  The theorem attacks an
architecture, not all solvers.

The smallest tested shape failure was $d=2$: the four terminal residuals had
70 bits versus a 34.7-bit $M^{2/3}$ target.  At $d=3$, both coordinates failed
(eight residuals versus $M^{1/3}=7.55$, each 146 bits versus
$M^{2/3}=56.97$).  At $d=8$, the exact counts were 256 distinct residuals of
8,693 bits from a 34,672-bit parent; the nominal target was 32.61 children of
1,063.29 bits.

\section{Amortization across branches}

The recovered QLEN1 inequality is correct in its explicit query-writing
model.  If $A(s)\le s^{k+o(1)}$, total written query length is
$V=\sum_i s_i$, and $S_{\max}=\max_i s_i$, then
\[
 \sum_i A(s_i)\le V S_{\max}^{k-1+o(1)}.
\]

\begin{lemma}[QLEN equality control]
For $q$ equal children of length $L$, the displayed bound is tight for the
padded cost function $A(s)=s^k$:
\[
 qL^k=(qL)L^{k-1}.
\]
At $q=M^{1/3}$ and $L=M^{2/3}$, its exponent is
\[
 \frac13+\frac{2k}{3}=1+\frac23(k-1).
\]
\end{lemma}

Thus query-length amortization faithfully charges the fanout; it does not
erase it.  When the recursive cost itself is $T(L)=L^t$, the exponent is
$1/3+(2/3)t$, whose fixed point is one.  That is the desired recursive
self-improvement \emph{after} an exact cover exists.  It does not construct
the cover or jointly answer branch-specific boundary questions.

The pointer-chain lemma shows why there is no representation-free stronger
rule: black-box branch work can be a true direct sum.  Conversely, the traffic
lemma shows that source scanning and output writing can be fully amortized at
the equality frontier.  Therefore the live question is semantic sharing, not
whether amortization is categorically possible.

\section{Typed cycles require two independent checks}

For a two-type cycle with active affine equations
\[
 C_A=w_{AB}+\lambda_{AB}C_B,
 \qquad C_B=w_{BA}+\lambda_{BA}C_A,
\]
and $\Lambda=\lambda_{AB}\lambda_{BA}<1$, the least fixed point is
\[
 C_A=\frac{w_{AB}+\lambda_{AB}w_{BA}}{1-\Lambda},
 \qquad C_B=w_{BA}+\lambda_{BA}C_A.
\]

\begin{proposition}[Contraction is not endpoint entry]
Cycle product below one removes dependence on an unknown starting exponent,
but it does not imply $C_A+S_A<\gamma$, where
$\gamma=2\cos(\pi/7)$.  For
\[
 \lambda_{AB}=1,\quad \lambda_{BA}=1/2,\quad
 w_{AB}=4/5,\quad w_{BA}=1/10,
\]
the product is $1/2$ but $C_A=9/5$.  With space floor $S_A=1/2$, the sum is
$23/10>\gamma$.
\end{proposition}

The executable also checked a passing control with the same scale product but
$w_{AB}=1/5,w_{BA}=1/10$: then $C_A=3/5$ and
$C_A+S_A=11/10<\gamma$.  Hence a typed construction needs both:
\begin{enumerate}
\item honest physical-scale contraction on every complete return cycle; and
\item the full least Bellman budgets below the endpoint.
\end{enumerate}

The delta-handle lemma supplies a third guard: size measures must be coherent
across types.  A cycle that returns to the same $\Theta(M)$ physical object
cannot be declared contracting merely by renaming it an $M^{2/3}$ handle and
omitting the $3/2$ expansion edge.

\section{A-code-dependent portfolios}

Code dependence is permitted by the intended contradiction route, but it
does not license uncharged dynamic information.

\begin{proposition}[No universal exponent extractor; conditional on a base
decider]
Assume there exists one polynomial-time exact SAT decider $B$.  There is no
total computable function $E$ that, from every code $e$ of a polynomial-time
exact SAT decider, returns a valid running-time exponent for that code.
\end{proposition}

\begin{proof}[Proof sketch]
Use the standard recursion theorem to obtain a code $e$ for a machine that
runs $B$ and then performs $n^{E(e)+1}$ irrelevant steps.  Since $E(e)$ is a
fixed integer, this machine remains a polynomial-time exact SAT decider, but
its running time exceeds the exponent $E(e)$ that was claimed to bound it.
\end{proof}

This does not obstruct the existing exponent-elimination recurrence, which
was designed not to know $K$.  It does obstruct any new portfolio argument
that silently obtains a valid $K$ from arbitrary source code or chooses caps
using such an extractor.

The padded family $A_d$ gives the operational version.  A portfolio member
defined by reproducing $A$'s parent trace is not acquisition-covariant under
semantics-preserving padding.  A viable $A$-dependent member must instead do
one of the following:
\begin{itemize}
\item statically normalize every allowed padding/interpreter wrapper with a
proved resource transport theorem;
\item extract a certificate before the padded work without already deciding
the parent; or
\item use $A$ only on genuinely smaller typed inputs and charge every access.
\end{itemize}

Finite portfolios remain valid: safely failing members may be tried in a
fixed order.  The missing statement is still worst-case coverage by at least
one favorably charged member, not a selector predicting which familiar
structural class applies.

\section{Weakest formulation not contradicted by this audit}

The following target survives every control above.

\begin{center}
\fbox{\begin{minipage}{0.92\textwidth}
\textbf{STREAMED-TYPED-SUCCESSOR2 (open).}  From the finite code of an assumed
uniform exact polynomial-time SAT machine $A$, construct a fixed finite typed
system.  On every ordinary CNF $F$ of physical size $M$, the system returns a
certified VALUE or a well-founded exact Boolean plan.  Children may be
ordinary explicit CNFs or source-queryable virtual objects.  For every
virtual object, declare its expanded logical length, handle length, exact
bit-decoder, aggregate decoder traffic, retained source state, and recursive
evaluator type.  Charge one shared acquisition, every failed member, every
source read, every decoder call, combination, replay, and all live state.
Measure all type edges against one coherent physical-size ledger.  Require
every reachable return cycle to contract honestly and require the least root
Bellman budgets to satisfy
\[
 C_0+S_0<2\cos(\pi/7).
\]
No parent-scale call to $A$, parent trace, undeclared succinct-SAT oracle, or
handle-to-expanded-size reset is allowed.
\end{minipage}}
\end{center}

This formulation intentionally permits $M^{1+o(1)}$ shared source work.  At
the clean explicit target it also permits $M^{1+o(1)}$ aggregate child
traffic and $M^{2/3+o(1)}$ sequential child space.  It does not require every
member to use the same syntax, nor does it require a witness-set cover when a
decision-only VALUE certificate is available.

\begin{longtable}{L{0.18\textwidth}L{0.19\textwidth}L{0.55\textwidth}}
\toprule
Relaxation & Audit result & Weakest surviving requirement\\
\midrule
\endfirsthead
\toprule
Relaxation & Audit result & Weakest surviving requirement\\
\midrule
\endhead
\bottomrule
\endfoot
Shared BatchAcquire & Not refuted & One semantic object with exact
counterfactual coverage, aggregate access, recursive closure, and charged
near-linear construction; not merely shared storage.\\
Compressed deltas & Handle-size shortcut refuted & Typed decoder or ordinary
materialization; count expanded logical scale and every source dependency.\\
Clause-first DAG & Universal equality-merge claim refuted & Partial member
with a proved progress potential or a charged semantic fallback covering the
hostile family.\\
Branch amortization & Traffic feasible; semantic sharing not automatic & A
constructive CNF-specific sharing theorem stronger than QLEN1, or accept the
direct-sum charge.\\
Typed cycles & Product-only acceptance refuted & Coherent physical scale plus
least Bellman time-space endpoint.\\
$A$-code portfolio & Trace/exponent shortcut refuted & Acquisition-covariant
static certificate or smaller-input uses of $A$; all failures charged.\\
\end{longtable}

\section{Executable results and reproducibility}

The bounded runner completed 1,020 declared checks in 2.38 seconds under a
25-second internal alarm and a 30-second shell cap.  It used random seed
220260108 only for disjoint pointer labels and terminal bits; every asserted
count is independently determined by the declared model.

\begin{center}
\small
\begin{tabular}{@{}L{0.32\textwidth}rL{0.46\textwidth}@{}}
\toprule
Section & Checks & Strongest finite observation\\
\midrule
Equality frontier & 16 & $qL=M$ exactly at tested powers of two; repeated
root scans cost $qM$.\\
Pivotal source reads & 264 & Every clause of $U_n$, $2\le n\le12$, was
deletion- and same-width-replacement-pivotal.\\
Payload/delta & 46 & Eight 48-bit handles named the same 2,280-bit payload
from a 2,329-bit parent.\\
Clause-first DAG & 592 & At $d=8$: 256 satisfiable distinct residuals, each
8,693 bits, from a 34,672-bit source.\\
QLEN equality & 32 & Direct $qL^k$ cost equaled the QLEN envelope in all 16
parameter rows.\\
Pointer-chain oracle & 48 & Exact $q(R+1)$ query count on all 16 controls.\\
Typed cycles & 9 & Same product $1/2$ produced one endpoint failure and one
pass; delta expansion product was one.\\
Padded-code trace & 13 & Full-trace exponent tracked arbitrary delays through
$d=12$.\\
\bottomrule
\end{tabular}
\end{center}

Executable:
\path{pnp_spines/final_link2_redteam.py}

Raw results:
\path{pnp_spines/final_link2_redteam_results.json}

SHA-256 of executable:
\path{761f441cf8820f88809c9563c204a6427b096bec3b993c531f802f0fe5c316d5}

SHA-256 of raw results:
\path{2d5f21e490de4a621d33f73112b04b77fd775319657a3a6a4e1005a1d51d32f0}

The report was informed by the recovered FINAL-LINK1 synthesis,
\path{pnp_spines/main8.tex}; the ARC ledger,
\path{research/agents/lead/2026-10-08_0032_LEAD_FINAL-LINK-ARC.tex}; the corpus
audit,
\path{research/agents/research/2026-10-08_003226_FINAL-LINK-CORPUS.tex}; and the
prior red-team report,
\path{research/agents/red/2026-10-07_1724_FINAL_LINK_REDTEAM.tex}.  It edits
none of those manuscripts.

\section{Bottom line}

The cleanest negative conclusion is architectural: delta handles,
clause-first paths, exact DAG merging, generic amortization, typed-cycle
products, and code access each omit a different charge when used alone.  The
cleanest positive conclusion is equally important: neither a linear source
scan nor linear aggregate child volume blocks the desired Williams region.
Searching for a stronger source-read lower bound is therefore not the best
next step unless it proves unavoidable repeated scans or a space-sensitive
tradeoff.

The highest-value remaining attempt is a concrete shared $H(A,F)$ with a
one-pass or few-pass Acquire phase, streamed child access, an exact
counterfactual interface, and a coherent typed Bellman certificate.  It must
be tested first on the payload family $C_w\wedge P$, the tagged address family
$D_d$, disjoint pointer-chain oracles, padded/interpreted $A_d$, and the
existing Tseitin/PHP controls.  Passing those controls would still be only a
candidate; proving all-input CNF coverage and the endpoint ledger is the
missing theorem.

% END PNP RESEARCH REPORT::FINAL-LINK2-REDTEAM::2026-10-08T01:02Z
\end{document}
